% GENERATED, edit the French unit or the translation table
% GENERATED from content/data/enrichment.json by mxbuild.enrichment; do not edit.
\mxunit{enrichment-invariants-solutions}{Parity and invariants}
\mxsubtitle{High school · reasoning \textperiodcentered{} Worked solutions \textperiodcentered{} 01}

\begin{exercise}
Start at 7 and add or subtract 4 each move. Can you reach 20? What about 23?
\end{exercise}
\begin{solution}
Every result is congruent to 7 modulo 4, hence to 3. Since 20 is congruent to 0, it is unreachable. Four additions of 4 give 23.
\end{solution}

\begin{exercise}
Five numbers, each 1 or $-1$, have product $-1$. Prove that changing two signs per move cannot make all five positive.
\end{exercise}
\begin{solution}
Changing two signs multiplies the product by $(-1)^{2} = 1$. The product stays $-1$, whereas five positive signs give 1.
\end{solution}

\begin{exercise}
Replace two integers a and b by $a + b$. What invariant does this preserve?
\end{exercise}
\begin{solution}
The sum of all integers stays constant: the two terms contributing $a + b$ are replaced by that same number.
\end{solution}

\begin{exercise}
Why need two states of the same parity not be connected by allowed moves? Give an example.
\end{exercise}
\begin{solution}
Parity gives only a necessary condition. With steps of 6 from 0, the target 2 is even but is not a multiple of 6.
\end{solution}

\begin{problem}[Investigate, justify, conclude]
An $8 \times  8$ chessboard is coloured alternately black and white. Remove two diagonally opposite corners. Can the 62 remaining squares be tiled by 31 dominoes of size $1 \times  2$, each covering two edge-adjacent squares? Count the removed colours and those covered by a domino. Compare with removing two neighbouring corners.
\end{problem}
\begin{solution}
Opposite corners have the same colour, leaving 30 squares of that colour and 32 of the other. Each domino covers one of each; 31 dominoes would need 31 of each, so tiling is impossible. Neighbouring corners have different colours. A tiling then exists: after removing the two top corners, place three horizontal dominoes in the first row and four in each of the other seven rows. Rotate this construction for the other sides.
\end{solution}

